% Chapter 2: Literature Review

\section{Introduction}

This chapter reviews existing research and systems related to livestock management, IoT in agriculture, data governance frameworks, and multi-agent systems. The review identifies research gaps and establishes the theoretical foundation for the Lesaka AI project.

\section{IoT in Agriculture}

\subsection{Smart Farming Systems}

Smart farming has emerged as a transformative approach to agriculture, leveraging IoT technologies to enhance productivity and sustainability. According to \cite{wolfert2017big}, IoT-enabled agricultural systems can reduce resource consumption by up to 30\% while increasing crop yields through precision monitoring.

\subsection{Livestock Monitoring}

Several studies have explored IoT-based livestock monitoring systems. \cite{riquelme2018smonitor} developed a wireless sensor network for cattle monitoring, focusing on GPS tracking and basic health parameters. However, their system lacked intelligent analysis capabilities and robust data governance mechanisms.

\subsection{Current Limitations}

Existing livestock monitoring systems face several limitations:

\begin{itemize}
    \item Limited data validation and integrity checks
    \item Lack of intelligent diagnostic capabilities
    \item Insufficient privacy protection for farmer data
    \item Poor integration with decision support systems
\end{itemize}

\section{Data Governance in Agricultural Systems}

\subsection{Database Normalization}

Third Normal Form (3NF) is a database design principle that eliminates redundancy and ensures data integrity. \cite{date2003database} emphasizes the importance of 3NF in systems requiring high data fidelity, such as agricultural telemetry systems.

\subsection{Privacy-by-Design}

Privacy-by-Design principles, as outlined by \cite{cavoukian2011privacy}, advocate for embedding privacy protections into system architecture from the outset. Agricultural systems must balance data utility with farmer privacy concerns.

\subsection{Role-Based Access Control}

RBAC systems provide fine-grained access control based on user roles \cite{sandhu1996role}. In agricultural contexts, different stakeholders (farmers, veterinarians, brokers) require varying levels of data access.

\section{Multi-Agent Systems}

\subsection{Agent Architectures}

Multi-agent systems (MAS) have been successfully applied in various domains including healthcare \cite{wooldridge2009introduction} and manufacturing \cite{jennings1998agent}. These systems enable distributed decision-making and autonomous problem-solving.

\subsection{Graph-Based Orchestration}

Directed Cyclic Graphs (DCG) provide a framework for modeling complex agent interactions \cite{diest2010graph}. Graph-based orchestration enables flexible routing of tasks among specialized agents based on contextual factors.

\subsection{Self-Healing Systems}

Self-healing systems can automatically detect and recover from failures \cite{ghosh2007self}. In agricultural IoT contexts, self-healing capabilities are crucial for maintaining system reliability in network-constrained environments.

\section{Comparison of Existing Systems}

\begin{table}[H]
\centering
\begin{tabular}{|l|l|l|l|l|}
\hline
\textbf{System} & \textbf{3NF} & \textbf{RBAC} & \textbf{MAS} & \textbf{Self-Healing} \\
\hline
System A \cite{ref1} & No & Limited & No & No \\
System B \cite{ref2} & Yes & No & Yes & Limited \\
System C \cite{ref3} & Limited & Yes & No & No \\
\textbf{Lesaka AI} & \textbf{Yes} & \textbf{Yes} & \textbf{Yes} & \textbf{Yes} \\
\hline
\end{tabular}
\caption{Comparison of Existing Livestock Management Systems}
\end{table}

\section{Research Gaps}

Based on the literature review, the following research gaps have been identified:

\begin{enumerate}
    \item Limited integration of 3NF database validation with intelligent agent systems
    \item Insufficient focus on privacy-by-design in agricultural IoT systems
    \item Lack of self-healing capabilities in rural network-constrained environments
    \item Minimal emphasis on Botswana-specific agricultural contexts
\end{enumerate}

\section{Conclusion}

The literature review reveals that while individual components of the proposed system have been explored in isolation, there is a lack of integrated solutions that combine robust data governance, intelligent agent orchestration, and self-healing capabilities specifically designed for Botswana's agricultural sector. The Lesaka AI project addresses these gaps through its comprehensive approach.

% Add references here - this is a placeholder
\bibliographystyle{apalike}
\bibliography{references}
